\documentclass[12pt]{article}

%% Basic document formatting
\usepackage{amsmath, amsthm, amssymb, amsfonts}
\usepackage{mathtools}
\usepackage{xspace}
\usepackage{thmtools}
\usepackage{graphicx}
\usepackage{setspace}
\usepackage{fancyhdr}
\usepackage{titling}
\usepackage[left=0.4in,right=0.4in,top=1in,bottom=1in]{geometry}
\usepackage{float}
\usepackage{hyperref}
\usepackage{mathrsfs}
\usepackage{tabularx}
\usepackage[utf8]{inputenc}
\usepackage[english]{babel}
\usepackage{framed}
\usepackage[dvipsnames]{xcolor}
\usepackage{environ}
\usepackage{tcolorbox}
\tcbuselibrary{theorems,skins,breakable}

\usepackage{enumitem}
\setlist[enumerate]{leftmargin=*}
\setlist[enumerate,1]{labelindent=\parindent}
\setlist[enumerate,2]{labelindent=0pt}

% Blackboard Bold
\newcommand{\Z}{\mathbb{Z}}                 % integers
\newcommand{\Zpl}{\mathbb{Z}_{+}}           % positive integers
\newcommand{\Znn}{\mathbb{Z}_{\geq 0}}      % nonnegative integers. 
\newcommand{\Q}{\mathbb{Q}}                 % rationals
\newcommand{\Qpl}{\mathbb{Q}_{+}}           % positive Rationals
\newcommand{\R}{\mathbb{R}}                 % reals
\newcommand{\Rpl}{\mathbb{R}_{+}}           % positive Reals
\newcommand{\C}{\mathbb{C}}                 % complex numbers
\newcommand{\F}{\mathbb{F}}                 % field

% Words
\newcommand{\st}{\text{ such that }}
\newcommand{\wrt}{\text{ with respect to }}
\newcommand{\with}{\text{ with }}
\newcommand{\ie}{\text{, i.e. }}
\newcommand*{\ditto}{\text{---}\texttt{"}\text{---}}

% Operators
\newcommand{\abs}[1]{\left|#1\right|}                   % absolute value
\newcommand{\norm}[1]{\abs{\abs{#1}}}
\newcommand{\floor}[1]{\left\lfloor #1 \right\rfloor}   % floor
\newcommand{\ceil}[1]{\left\lceil #1 \right\rceil}      % ceiling
\newcommand{\powset}[1]{\mathscr{P}(#1)}

% Category Theory 
\renewcommand{\hom}{\text{Hom}}
\newcommand{\homspace}[3]{\hom_{#1}(#2, #3)}
\newcommand{\coproduct}[9]{
  \begin{tikzcd}[ampersand replacement=\&]
      \& \& #1 \arrow[dll, bend right, "#6"] \arrow[dl, "#5"] \\
   #4 \& #3 \arrow[l, "#9"] \\
      \& \& #2 \arrow[ull, bend left, "#8"] \arrow[ul, "#7"]
  \end{tikzcd}
}

\newcommand{\product}[9]{ 
  \begin{tikzcd}[ampersand replacement=\&]
      \& \& #3 \\
      #1 \arrow[r, "#7"] \arrow[urr, bend left, "#5"] \arrow[drr, bend right, "#6"] \& #2 \arrow[ur, "#8"] \arrow[dr, "#9"] \\ 
      \& \& #4
  \end{tikzcd}
}


% Algebra 
\newcommand{\Syl}[2]{\text{Syl}_{#1}{(#2)}}           % set of Sylow-p subgroups 
\newcommand{\<}{\langle}                % \<x\>, subgroup generated by x 
\renewcommand{\>}{\rangle}
\renewcommand{\index}[2]{[#1 : #2]}
\newcommand{\id}{\text{id}}             % identity element
\newcommand{\lhdneq}{%
  \mathrel{\ooalign{$\lneq$\cr\raise.22ex\hbox{$\lhd$}\cr}}}
\newcommand{\order}[1]{\text{o}(#1)}    % order of an element
\newcommand{\spec}{\operatorname{Spec}}
\newcommand{\rad}{\operatorname{rad}}   % radical of an ideal 
\newcommand{\sym}[1]{\mathfrak{S}_{#1}}
\newcommand{\aut}[1]{\text{Aut}(#1)} 
\newcommand{\gal}[2]{\text{Gal}(#1 / #2)} 
\newcommand{\inn}[1]{\text{Inn}(#1)} 
\let\oldcong\cong
\let\oldequiv\equiv
\renewcommand{\cong}{\oldequiv}
\renewcommand{\equiv}{\oldcong}

\newcommand{\triangleleftneq}{\mathrel{\ooalign{%
  $\trianglelefteq$\cr
  \hidewidth\raise0.06ex\hbox{$\not\mkern4mu$}\hidewidth\cr
}}}

\makeatletter % cycle
\newcommand{\cyc}[1]{(\mathbf{\cyc@process#1\relax})}
\def\cyc@process#1#2\relax{%
	#1%
	\ifx\relax#2\relax
	\else
		\,\cyc@process#2\relax
	\fi
}
\makeatother

\newcommand{\GL}[2]{\text{GL}_{#1}(#2)}                             % general linear group
\newcommand{\SL}[2]{\text{SL}_{#1}(#2)}                             % special linear group
\newcommand{\gl}[2]{\mathfrak{gl}_{#1}(#2)}
\renewcommand{\sl}[2]{\mathfrak{sl}_{#1}(#2)}
\newcommand{\so}[2]{\mathfrak{so}_{#1}(#2)}
\newcommand{\su}[1]{\mathfrak{su}(#1)}

% Linear Algebra
\newcommand{\bmat}[1]{\begin{bmatrix}#1\end{bmatrix}}   % bracketed matrix
\newcommand{\rank}{\operatorname{rank}}                 % rank
\newcommand{\nullity}{\operatorname{nullity}}           % nullity
\newcommand{\tr}{\operatorname{tr}}

% Combinatorics
\newcommand{\qnum}[1]{\left[#1\right]_q}                % q-analog
\newcommand{\qfact}[1]{\left[#1\right]!_q}              % q-analog factorial 
\newcommand{\qbinom}[2]{\binom{#1}{#2}_q}               % Gaussian binomial coefficient

% Topology/Analysis
\newcommand{\ball}[2]{\text{B}_{#1}(#2)}  % B_r(x): open r-balls around x
\newcommand{\diam}{\text{diam}}           % diamter of a set in metric space 
\newcommand{\lowint}[2]{\underline{\int_{#1}^{#2}}}
\newcommand{\upint}[2]{\overline{\int}_{#1}^{#2}}

% Misc Notation
\newcommand{\oneton}[1]{\{1, \ldots, #1\}}
\newcommand{\defeq}{\vcentcolon=}   % :=
\newcommand{\eqdef}{=\vcentcolon}   % =: 
\renewcommand{\bf}[1]{\textbf{#1}}
\newcommand{\im}{\operatorname{im}}
\newcommand{\fnrest}[2]{\left.#1\right|_{#2}}
\newcommand{\isom}{\overset{\sim}{\rightarrow}} 


\renewcommand{\thesection}{\arabic{section}.}
\renewcommand{\thesubsection}{(\alph{subsection})}

\newtheorem{claim}{Claim}
\newtheorem{theorem}{Theorem}
\newtheorem{proposition}{Proposition}
\newtheorem*{lemma}{Lemma}

%% Headers & title setup
\newcommand{\course}{Math140C}
\newcommand{\myname}{Jay Ser}
\setlength{\headheight}{14.5pt}
\pagestyle{fancy}
\fancyhf{}
\renewcommand{\headrulewidth}{0.4pt}
\lhead{\course}
\rhead{\myname}
\cfoot{\thepage}
\setlength{\droptitle}{-4em} 
\title{\course\ - HW \#4}
\author{\myname}
\date{2026.05.10} 

\begin{document}
\maketitle
\thispagestyle{fancy}

%------------------------------------------------------------------------------%
\newcommand{\M}{\mathcal{M}}
\newcommand{\scrC}{\mathscr{C}}
\newcommand{\sigalg}[1]{\sigma\text{-alg}(#1)}

\section{}
Let $\M$ be an algebra on a set $X$. 
\subsection{} 
Show $\emptyset, X \in \M$ 

\noindent \dotfill

For any $A \in \M$, $A \cup (X \setminus A) = X \in \M$. 
And hence $X \setminus X = \emptyset \in \M$. 

\subsection{} 
Show that if $A_1, \ldots, A_n \in \M$, then $\bigcap_{i = 1}^{n} A_i$. 
Similarly, supposing $\M$ is a $\sigma$-algebra, show that if $\{A_n\}_{n \in \Zpl} \subset \M$, then $\bigcap_{i = 1}^{\infty} A_i \in \M$.

\noindent\dotfill

For the first part, it clearly suffices to show that if $A, B \in \M$, then $A \cap B \in \M$. 
Well, deMorgan's law says  
$$X \setminus ((X \setminus A) \cup (X \setminus B)) = A \cap B,$$
and the left hand side is in $\M$, which shows $A \cap B \in \M$.  

For $\sigma$-algebras, once again, use deMorgan's Law: 
$$X \setminus \bigcup_{i = 1}^{\infty} (X \setminus A_i) = \bigcap_{i = 1}^{\infty} A_i$$
If $\{A_i\}_{i \in \Zpl} \subset \M$, then the left hand side is in $\M$, so $\bigcap_{i} A_i \in \M$. 

\section{} 
\subsection{} 
Let $I$ be any set and $\M_i$ a $\sigma$-algebra on $X$ for every $i \in I$. 
Prove $\bigcap_{i \in I} \M_i$ is a $\sigma$-algebra on $X$. 

\noindent\dotfill 

Let $\M = \bigcap_{i \in I} \M_i$. 
Take $A \in \M$. 
For each $i \in I$, $A \in \M_i$, so it follows $X \setminus A \in \M_i$. 
This holds for each $i \in I$, so $X \setminus A \in \M$. 
Closure under arbitrary unions holds by the same line of reasoning. 

\subsection{} 
Take any $\mathscr{C} \subset \mathscr{P}(X)$. 
Show there is a unique smallest $\sigma$-algebra $\M$ on $X$ that contains $\mathscr{C}$. 

\noindent\dotfill

Let $\{\M_i\}_{i \in I}$ be the collection of all $\sigma$-algebras that contain $\mathscr{C}$. 
Then 
$$\bigcap_{i \in I}\M_i$$
is clearly the unique smallest $\sigma$-algebra that contains $\mathscr{C}$. 

\section{} 
\subsection{} 
Let $\scrC_1, \scrC_2 \subset \mathscr{P}(X)$. 
Show $\scrC_1 \subset \sigalg{\scrC_2} \iff \sigalg{\scrC_1} \subset \sigalg{\scrC_2}$ 

\noindent\dotfill 

The reverse direction is clear since $\scrC_1 \subset \sigalg{\scrC_1}$. 
The forward direction is also immediate because $\sigalg{\scrC_1}$ is by definition the smallest $\sigma$-algebra containing $\scrC_1$, hence $\scrC_1 \subset \sigalg{\scrC_2}$ implies $\sigalg{\scrC_1} \subset \sigalg{\scrC_2}$ by definition. 

\subsection{} 
Let $\mathcal{I} = \{(a, b) \subset \R \mid a < b \text{ and } a, b \in \Q\}$. 
Explain why $\mathcal{B}_{\R} = \sigalg{\mathcal{I}}$. 

\noindent\dotfill 

$\mathcal{I}$ is a subset of all the open sets of $\R$, so $\sigalg{\mathcal{I}} \subset \mathcal{B}_{\R}$ is immediate. 
For the reverse inclusion, observe that open subsets in $\R$ form a subset of $\sigalg{\mathcal{I}}$: 
suppose $U \subset \R$ is open and pick any $x \in U$. 
By definition, there exists $r > 0$ such that $\ball{r}{x} \subset U$. 
Since $\Q$ is dense in $\R$, one can choose $a, b \in \Q$ such that $(a, b) \subset \ball{r}{x} \subset U$. 
Hence $U$ is a union intervals with rational endpoints, i.e., 
$$U \in \sigalg{\mathcal{I}}.$$
Since $U$ is an arbitrary open set in $\R$, the above shows the set of all open sets in $\R$ is a subset of $\sigalg{\mathcal{I}}$. 
$\mathcal{B}_{\R} \subset \sigalg{\mathcal{I}}$ follows by part (a). 

If one knows a bit of topology, this problem is immediate from the fact that open intervals with rational endpoints form a basis for the regular topology of $\R$. 

\section{}
Let $\scrC = \{\{x\} \mid x \in \R\}$ be the collection of singleton subsets of $\R$. 
Find a characterization for the sets belonging to $\sigalg{\scrC}$. 
Specifically, find an explicit condition, and show that the collection of sets satisfying your condition 
\begin{itemize} 
  \item includes every set in $\scrC$
  \item is a $\sigma$-algebra 
  \item is contained in $\sigalg{\scrC}$. 
\end{itemize} 

\noindent\dotfill 

Note that the outlined conditions indeed show such a collection is $\sigalg{\scrC}$. 
Well, I claim that $\powset{\R} = \sigalg{\scrC}$. 
It is trivially true that $\powset{\R}$ contains all singletons in $\R$, i.e., every set in $\scrC$, and is a $\sigma$-algebra. 
To see that $\powset{\R} \subset \sigalg{\scrC}$, take any $S \subset \R$. 
By the definition of a set, 
$$S = \bigcup_{x \in S} \{x\},$$
hence $S \in \sigalg{\scrC}$. 
This proves $\powset{R} = \sigalg{\scrC}$. 

\section{} 
Let $R = \prod_{i = 1}^{p} I_i \in \mathcal{R}$ where each $I_i \subset \R$ is connected. 
A coordinate-wise partition of $R$ is a partition $P$ of $R$ such that there exists finite partitions $Q_i$ of $I_i$ for each $i \in \oneton{p}$ such that 
\begin{itemize}
  \item for each $i \in \oneton{p}$, each set in $Q_i$ is connected 
  \item $P = \{\prod_{i = 1}^{p} J_i \mid J_i \in Q_i \text{ for each } i \in \oneton{p}\}$ 
\end{itemize}

\subsection{} 
Show that if $P$ is a coordinate-wise partition, then 
\begin{equation} \label{eq:5_a_wts}
m(R) = \sum_{D \in P} m(D) = \sum_{J_1 \in Q_1} \sum_{J_2 \in Q_2} \cdots \sum_{J_p \in Q_p} m(\prod_{i = 1}^{p} J_i).
\end{equation}

\noindent\dotfill

After observing that each element of $P$ is itself a bounded rectangle, the second equality in \eqref{eq:5_a_wts} holds by definition, so it suffices to show the equality between the first and third expressions in \eqref{eq:5_a_wts}. 
Let $a_i, b_i$ be for each $i \in \oneton{p}$ such that 
\begin{gather*}
  \prod_{i = 1}^{p} (a_i, b_i) \subset R \subset \prod_{i = 1}^{p} [a_i, b_i] \\ 
  m(R) = \prod_{i = 1}^{p} (b_i - a_i)
\end{gather*}

Consider any index $i \in \oneton{p}$. 
Without loss of generality, suppose $I_i = [a_i, b_i]$ 
(in the treatment below, I will write all segments as closed, just for notational simplicity; 
the distinction of a segment being closed is unimportant for part (a)). 
Let $Q_i$ be as defined in the problem statement. 
By the restrictions on $Q_i$, there exist $x_{i0}, x_{i1}, \ldots, x_{it_i}$ for some integer $t_i$ such that $x_{i0} = a_i$, $x_{it_i} = b_i$, and 
$$Q_i = \{[x_{i0}, x_{i1}], [x_{i1}, x_{i2}], \ldots, [x_{i(t_i - 1)}, x_{it_i}]\}.$$
By definition of the $x_{ij}$'s,  
$$b_i - a_i = \sum_{j = 1}^{t_i} (x_{ij} - x_{i(j - 1)}).$$
With $t_i$ and $x_{ij}$ defined as above for all $i \in \oneton{p}$ and $j \in \oneton{t_i}$, 
\begin{align}
  m(R) = \prod_{i = 1}^{p} (b_i - a_i) 
  & = \prod_{i = 1}^{p} (\sum_{j = 1}^{t_i} (x_{ij} - x_{i(j-1)})) \nonumber \\ 
  & = \sum_{j_i = 1}^{t_1}\sum_{j_2 = 1}^{t_2} \cdots \sum_{j_p = 1}^{t_p} \prod_{k = 1}^{p} (x_{kj_i} - x_{k(j_i - 1)}) \label{eq:5_a_dist} \\ 
  & = \sum_{J_1 \in Q_1}\sum_{J_2 \in Q_2} \cdots \sum_{J_p \in Q_p} m(\prod_{k = 1}^{p} J_k), \nonumber
\end{align}
where \eqref{eq:5_a_dist} is due to the distributive property.

\subsection{} 
Show that if $R = \bigsqcup_{k = 1}^{n} A_k$ for disjoint squares $A_1, \ldots, A_n \in \mathcal{R}$, then for every $k \in \oneton{n}$, there is a coordinate-wise partition $P_k$ of $A_k$ such that $P = \bigcup_{k = 1}^{n} P_k$ is a coordinate-wise partition of $R$. 

\noindent\dotfill

Associate with each $A_k$ points $a_{ki}, b_{ki}$ for all $i \in \oneton{p}$ such that 
$$\prod_{i = 1}^{p}(a_{ki}, b_{ki}) \subset A_k \subset \prod_{i = 1}^{p} [a_{ki}, b_{ki}]$$
For each $i \in \oneton{p}$, define 
$$T_i = \{a_{ki} \mid k \in \oneton{n}\} \cup \{b_{ki} \mid k \in \oneton{n}\}$$ 
Clearly, $T_i$ is a finite collection of points lying in the segment $[a_i, b_i]$. 
Ordering the points of $T_i$ under $<$ as in $T_i = \{p_0, p_1, \ldots, p_{t_i}\}$ for some $t_i$, each $p_{j} < p_{j + 1}$, let 
$$Q_i = \{(p_0, p_1), (p_1, p_2), \ldots, (p_{t_i - 2}, p_{t_i - 1}), (p_{t_i - 1}, p_{t_i})\} \cup \bigcup_{j = 0}^{t} \{p_{j}\}$$
By construction, the $Q_i$'s are a finite partition of the corresponding $I_i$'s such that 
$$P = \left\{\prod_{i = 1}^{p} J_i \mid J_i \in Q_i \text{ for each } i \in \oneton{p}\right\}$$
is a coordinate-wise partition of $R$. 
Crucially, $P$ induces a coordinate-wise partition for each $A_k$.
For a fixed dimension $i \in \oneton{p}$ and fixed rectangle $A_k$ for $k \in \oneton{n}$, one can form the segment $I_{ki}$ defining the $i$th dimension of $A_k$ by taking union of elements of $Q_i$; 
namely, that the endpoints of $I_{ki}$ are added to $Q_i$ as singletons means the elements of $Q_i$ can form open, closed, and half-open intervals with endpoints $a_{ki}, b_{ki}$. 
In other words, given a rectangle $A_k$, each $Q_i$ admits a subset partitioning $I_{ki}$ that defines the coordinate-wise partition for $A_k$. 
Furthermore, it is clear that the only shared points among the $P_k$'s are rectangles of measure zero; 
in each coordinate, the segments $I_{k_1 i}$ and $I_{k_2 i}$ for $k_1 \neq k_2 \in \oneton{n}$ share at most their endpoints.
Conversely, all elements of $P$ belong to a coordinate-wise partition of at least one of the $A_k$'s. 

Writing $P_k$ as this induced coordinate-wise partition of $A_k$, the conclusion obtained from the above paragraph is that $P = \bigcup_{k = 1}^{n} P_k$. 

\subsection{}
Let the notation be as in part (b). 
Combining the results from parts (a) and (b), 
\begin{align*}
  m(R) = \sum_{D \in P} m(R) 
  & = \sum_{k = 1}^{n} \sum_{D \in P_k} m(D) \\ 
  & = \sum_{k = 1}^{n} m(A_k),
\end{align*}
as was to be shown.

\section{}
\subsection{}
Let $R \in \mathcal{R}$ and let $R_1, \ldots, R_n \in \mathcal{R}$ be disjoint subsets of $R$. 
Show there is a coordinate-wise partition $P$ of $R$ so that each $R_k$ is a union of sets from $P$.
Conclude that $R \setminus (R_1 \cup \ldots \cup R_n)$ can be expressed as a finite union of disjoint rectangles.

\noindent\dotfill 

For the first part of this question, proceed more or less exactly as in Q5(b). 
The only thing to note is that excluding sets of $P$ of measure 0, each $R_k$ is a disjoint union of sets of $P$; 
since the $R_k$'s are themselves disjoint, the sets of $P$ used to build the rectangles $R_k$ contribute to only one $R_k$ if they have measure greater than 0.  
Since $P$ partitions $R$ into subrectangles, and $P$ induces further coordinate-wise partitions of the $R_k$'s, $R \setminus (R_1 \cup \ldots \cup R_n)$ can be expressed as a union of disjoint rectangles. 

\subsection{} 
Prove that every $A \in \mathscr{E}$ can be expressed as a finite union of disjoint rectangles. 

\noindent\dotfill 

By definition, $A = \bigcup_{k = 1}^{n} A_k$ for some integer $n$, $A_k \in \mathcal{R}$, where the $A_k$'s may not be pairwise disjoint. 
To show that any $A \in \mathcal{E}$ can be expressed as the union of disjoint rectangles, induct on $n$. 
The $n = 1$ is trivially true. 
Consider $n > 1$ with the proposition true for $n - 1$. 
Rewrite
$$A = \bigcup_{k = 1}^{n - 1} A_k \cup A_n$$
$A' \defeq \bigcup_{k = 1}^{n - 1}A_k$ is an element of $\mathscr{E}$, so by the inductive hypothesis, $A'$ can be expressed as the union of disjoint rectangles. 
If $A'$ is disjoint from $A_n$, then $A$ is the union of pairwise disjoint rectangles. 
Suppose $A'$ intersects $A_n$; 
particularly, there are some number of disjoint rectangles $B_1, \ldots, B_t$ of $A'$ that intersect $A_n$. 
Clearly, each $B_j \cap A_n$ is a rectangle subset of $A_n$, hence part (a) says that 
$$A_n \setminus (A_n \cap \bigcup_{j = 1}^{t} B_t)$$
can be written as a union of disjoint rectangles. 
In other words, part of $A_n$ that doesn't intersect $A'$ can be writ as a union of disjoint rectangles. 
Hence $A$ can be writ as a union of disjoint rectangles. 

\subsection{} 
Let $A \in \mathscr{E}$ and $R_1, \ldots, R_k \in \mathcal{R}$ be one set of pairwise disjoint rectangles, $S_1, \ldots, S_l \in \mathcal{R}$ another set of pairwise disjoint rectangles, and $A = \bigsqcup_{i = 1}^{k} R_i = \bigsqcup_{j = 1}^{k} S_j$. 
Show $\sum_{i = 1}^{k} m(R_i) = \sum_{j = 1}^{l} m(S_j)$. 

\noindent\dotfill 

For any $i \in \oneton{k}$, $R_i \subset A = \bigsqcup_{j = 1}^{l} S_j$ means 
\begin{align}
  R_i & = R_i \cap \bigsqcup_{j = 1}^{l} S_j \nonumber \\ 
      & = \bigsqcup_{j = 1}^{l} (R_i \cap S_j) \label{eq:6_c_sqcup}
\end{align}
(and indeed, simple set theory shows \eqref{eq:6_c_sqcup} is still a disjoint union). 
The intersection of rectangles is clearly another rectangle; 
\eqref{eq:6_c_sqcup} shows the rectangle $R_i$ is a disjoint union of rectangles. 
Hence Q5(c) applies to give 
$$m(R_i) = \sum_{j = 1}^{l} m(R_i \cap S_j)$$
Not only is the choice of $i$ arbitrary, but also a symmetric analysis can be to write each $S_j$ as a disjoint union of rectangles, yielding 
$$m(S_j) = \sum_{i = 1}^{k} m(S_j \cap R_i).$$
Putting these results together,  
\begin{align*}
  \sum_{i = 1}^{k} m(R_i) & = \sum_{i = 1}^{k} \sum_{j = 1}^{l} m(R_i \cap S_j) \\ 
                          & = \sum_{j = 1}^{l} \sum_{i = 1}^{k} m(R_i \cap S_j) = \sum_{j = 1}^{l} m(S_j), 
\end{align*}
as was to be shown. 

\end{document}
